% =============================================================================
%  Lab 8 - Next Instruction Address Generation II
% =============================================================================

\chapter{Lab 8 --- Next Instruction Address Generation II}

\section{Objective}
% Two or three sentences, your own words.

\section{Specification}

\begin{table}[htbp]
\centering
\caption{PC source selection (Type C). Only the relative path is conditional.}
\label{tab:lab08-sel}
\begin{tabular}{cl||ll}
\toprule
\# & \texttt{Sel} & \texttt{NextPC} & \\
\midrule
0 & \texttt{NIAG\_SEL\_NSA} & \texttt{NSA} & unconditional \\
1 & \texttt{NIAG\_SEL\_REL} & \texttt{Taken ? Rel : NSA} & \textbf{conditional} \\
2 & \texttt{NIAG\_SEL\_REG} & \texttt{RegAddr} & unconditional \\
3 & \texttt{NIAG\_SEL\_ABS} & \texttt{Abslt} & unconditional \\
\bottomrule
\end{tabular}
\end{table}

\begin{table}[htbp]
\centering
\caption{Branch conditions (Type E, exhaustive over the two ALU flags).}
\label{tab:lab08-cond}
\begin{tabular}{lll}
\toprule
code & meaning & Taken $=$ \\
\midrule
\texttt{NIAG\_COND\_NONE} & always   & $1$ \\
\texttt{NIAG\_COND\_EQ}   & $=0$     & $Z$ \\
\texttt{NIAG\_COND\_NEQ}  & $\neq 0$ & $\overline{Z}$ \\
\texttt{NIAG\_COND\_GTZ}  & $>0$     & $\overline{N}\cdot\overline{Z}$ \\
\texttt{NIAG\_COND\_GEZ}  & $\geq 0$ & $\overline{N}$ \\
\texttt{NIAG\_COND\_LTZ}  & $<0$     & $N$ \\
\texttt{NIAG\_COND\_LEZ}  & $\leq 0$ & $N + Z$ \\
\bottomrule
\end{tabular}
\end{table}

% REQUIRED, one or two sentences: three of these pairs are complements. Name
% them, and say what that tells you about how much hardware the seven
% conditions actually cost.

\subsection*{Absolute target formation}

\[
\texttt{Abslt} = \underbrace{\texttt{NSA}(31..28)}_{\text{4 inherited}} \;\&\;
\underbrace{\texttt{Index}(25..0)}_{\text{26 encoded}} \;\&\;
\underbrace{\texttt{"00"}}_{\text{2 always zero}}
\]

% REQUIRED: work through the real J from the Lab 12 program (0x14,
% 0x08000002) and show the target is 0x08.

\section{Design}
% Answer all three:
%   (a) Why are Sel and Cond two SEPARATE control fields rather than one? What
%       goes wrong if you gate the jump paths with Taken as well?
%   (b) Why does a not-taken branch fall back to NSA rather than needing its own
%       select code?
%   (c) The top four bits of a jump target are inherited, not encoded. State the
%       range limit that produces, explain where the number comes from, and say
%       what instruction you would use to go further.

\section{Implementation}

\vhdlfile{code/lab08/NIAG_Cond.vhd}{\texttt{NIAG\_Cond} --- the six branch-condition equations}

\vhdlfile{code/lab08/NIAG_MUX.vhd}{\texttt{NIAG\_MUX} --- the conditional gate and the four-way selection}

\vhdlfile{code/lab08/NIAG.vhd}{\texttt{NIAG} --- the complete address generator, six components}

\section{Verification}
% Vectors, then the coverage argument.
%
% Two different arguments apply here and you should say so:
%   * NIAG_Cond is Type E - two flags, seven conditions, 28 cases. Small enough
%     to enumerate exhaustively, so there is no argument to make: just do it.
%   * The full NIAG is Type C - test one representative of each PC source, plus
%     the taken/not-taken pair that distinguishes conditional from unconditional.

\begin{table}[htbp]
\centering
\caption{The scripted sequence.}
\label{tab:lab08-sequence}
\begin{tabular}{clllp{3.6cm}}
\toprule
\# & what happens & control & expected PC & what it proves \\
\midrule
1 & & & & \\
2 & & & & \\
3 & & & & \\
\bottomrule
\end{tabular}
\end{table}

\section{Results}
% Evidence: E1, E2, E4, E7.

\figplaceholder{the scripted-sequence waveform --- the PC walking through
sequential, not-taken branch, taken branch, \texttt{J}, register jump, stall and
backward branch, with EVERY transition annotated by which source it took}
% This is the centrepiece figure of the chapter. An unannotated 30-signal
% waveform proves nothing. Label each PC change with NSA / Rel / Abslt / RegAddr.

\figplaceholder{the RTL Viewer, showing the six components and how they connect}

% E4: register count. Expect 32 - the PC - and nothing else. Everything else in
%     this block is combinational. Say so explicitly.

\section{Discussion}
% REQUIRED. One bug, how you found it, what it cost.
%
% If you need a topic: did your GTZ condition fire when the result was zero?
% Would your original test vectors have caught it? What does that tell you about
% choosing test values for a condition code?
